\documentclass[11pt,a4paper]{article}
\usepackage{cursus}
\author{dr. Maarten Slembrouck}

\begin{document}

\labotitel{Labo 5 -- Datavisualisatie I}{Vertel het verhaal met plotly}{Cursustekst bij het notebook \code{lab5\_data\_visualization\_I.py}}

\section*{Waarover gaat dit labo?}

In twee labo's maakte je van rommelige exports nette, gecombineerde
tabellen. Maar een coach leest geen tabel met 111 rijen. Een goede
grafiek toont in één oogopslag wat in de tabel verborgen zit: wie het
snelst is, wie vooruitgaat, hoe zwaar een training was.

\textbf{Na dit labo kun je} een grafiektype kiezen dat past bij een
vraag, de standaardgrafieken maken met \code{plotly.express}, extra
informatie tonen met kleur, facetten en grootte, en grafieken
interactief maken in marimo.

Alle grafieken in deze tekst zijn gemaakt met de code die erboven staat,
op de data van de club. In het notebook zijn ze interactief: je kunt
inzoomen en met de muis over een punt gaan voor meer informatie.

\begin{python}[hide]
import pandas as pd
import plotly.express as px
import plotly.io as pio

pio.templates.default = "plotly_white"

def clean_positions(position_column):
    position_map = {
        "forward": "Forward", "fwd": "Forward",
        "midfielder": "Midfielder", "mid": "Midfielder",
        "defender": "Defender", "def": "Defender",
        "goalkeeper": "Goalkeeper", "gk": "Goalkeeper",
    }
    return position_column.str.strip().str.lower().map(position_map)

def fix_sprint_units(seconds_column):
    return seconds_column.where(seconds_column < 20, seconds_column / 1000)

REFERENCE_DATE = pd.Timestamp("2026-09-14")
athletes = pd.read_csv(f"{DATA_DIR}/athletes.csv").assign(
    position=lambda d: clean_positions(d["position"]),
    weight_kg=lambda d: d["weight_kg"].str.replace(",", ".").astype(float),
    age=lambda d: (REFERENCE_DATE - pd.to_datetime(d["birthdate"])).dt.days // 365,
).drop_duplicates(subset=["name", "birthdate"], keep="first")
tests = pd.read_csv(f"{DATA_DIR}/tests.csv").assign(
    sprint_10m_s=lambda d: fix_sprint_units(d["sprint_10m_s"]),
    sprint_30m_s=lambda d: fix_sprint_units(d["sprint_30m_s"]),
)
combined = athletes.merge(tests, on="athlete_id", how="inner")
hr_session = pd.read_csv(f"{DATA_DIR}/hr_session.csv")
hr_session["timestamp"] = pd.to_datetime(hr_session["date"]) + pd.to_timedelta(hr_session["time_s"], unit="s")
hr_with_age = hr_session.merge(athletes[["athlete_id", "age", "name"]], on="athlete_id", how="left")
hr_with_age["hr_max"] = 220 - hr_with_age["age"]
hr_with_age["pct_hr_max"] = 100 * hr_with_age["heart_rate_bpm"] / hr_with_age["hr_max"]
hr_zoned = hr_with_age.assign(hr_zone=lambda d: pd.cut(
    d["pct_hr_max"], bins=[0, 60, 70, 80, 90, 1000],
    labels=["Zone 1", "Zone 2", "Zone 3", "Zone 4", "Zone 5"]))
\end{python}

% ---------------------------------------------------------------------------
\section{De data voor dit labo}

Het notebook begint met één cel die de volledige opkuis van labo 3 en 4
herhaalt. Daarna heb je twee tabellen:

\begin{itemize}
  \item \code{combined}: één rij per atleet per testmoment, met positie,
        geslacht, gewicht, sprinttijden, spronghoogte en Yo-Yo-niveau;
  \item \code{hr\_zoned}: de training per seconde, met hartslag,
        snelheid, naam en hartslagzone.
\end{itemize}

\begin{python}
combined.shape, hr_zoned.shape
\end{python}

Beide zijn \textbf{lang}: één rij per meting. Dat is precies de vorm die
\code{plotly.express} verwacht (zie labo 4).

Daarnaast legt het notebook vooraf één kleur per positie vast, in een
dictionary. Zo is ``Forward'' in elke grafiek dezelfde kleur, en moet de
lezer de legende maar één keer leren.

\begin{python}[noout]
POSITION_COLORS = {
    "Forward": "#2a78d6",
    "Midfielder": "#eb6834",
    "Defender": "#1baf7a",
    "Goalkeeper": "#eda100",
}
\end{python}

Een kleur schrijf je als \emph{hexcode}: een hekje gevolgd door zes
tekens, twee voor rood, twee voor groen en twee voor blauw.

% ---------------------------------------------------------------------------
\section{Principes: de juiste grafiek voor de juiste vraag}

Elke grafiek beantwoordt een vraag. Begin daarom bij de vraag, niet bij
het grafiektype dat je toevallig kent.

\begin{center}
\rowcolors{2}{tint}{white}
\begin{tabularx}{\textwidth}{@{}L{2.8cm}Y L{3.3cm}@{}}
\toprule
\textbf{Soort vraag} & \textbf{Voorbeeld} & \textbf{Geschikte grafiek} \\
\midrule
Verdeling & Hoe zijn de sprinttijden gespreid? & histogram, boxplot \\
Vergelijking & Welke positie springt gemiddeld het hoogst? & staafdiagram, boxplot \\
Verband & Betekent een snellere sprint een hogere sprong? & spreidingsdiagram \\
Evolutie & Stijgt de hartslag van deze atleet? & lijngrafiek \\
\bottomrule
\end{tabularx}
\end{center}

\subsection{Eerlijke assen}

Een grafiek kan even goed misleiden als informeren. De bekendste valkuil
is een \textbf{afgeknotte as}. Hieronder twee keer dezelfde gemiddelden.
In de eerste grafiek begint de y-as bij 0:

\begin{python}[fig=eerlijk]
sprint_means = combined.groupby("position", as_index=False)["sprint_30m_s"].mean()
fig = px.bar(
    sprint_means, x="position", y="sprint_30m_s",
    color="position", color_discrete_map=POSITION_COLORS,
    title="Gemiddelde sprinttijd over 30 m per positie",
    labels={"sprint_30m_s": "Gemiddelde sprinttijd (s)", "position": "Positie"},
)
fig.update_yaxes(rangemode="tozero")
\end{python}

\code{as\_index=False} zorgt ervoor dat \code{position} een gewone kolom
blijft (en geen index wordt), zodat je ze als \code{x} kunt gebruiken.
In de tweede grafiek begint de y-as net onder de kleinste waarde:

\begin{python}[fig=misleidend]
fig = px.bar(
    sprint_means, x="position", y="sprint_30m_s",
    color="position", color_discrete_map=POSITION_COLORS,
    title="MISLEIDEND - gemiddelde sprinttijd over 30 m per positie",
    labels={"sprint_30m_s": "Gemiddelde sprinttijd (s)", "position": "Positie"},
)
fig.update_yaxes(range=[5.1, 5.55])
\end{python}

Dezelfde kleine verschillen (hooguit 0,3 seconden) lijken nu enorm. Een
staaf toont een grootte met haar \emph{lengte}; dan moet die lengte bij
nul beginnen. Om dezelfde reden gebruik je nooit twee verschillende
y-assen in één grafiek: de lezer kan de twee schalen niet vergelijken.

% ---------------------------------------------------------------------------
\section{Grafieken maken met \code{plotly.express}}

\code{plotly.express} is een bibliotheek om met weinig code interactieve
grafieken te maken. Je importeert ze als \code{px}:

\begin{python}[noout]
import plotly.express as px
\end{python}

Elke functie volgt hetzelfde patroon:

\begin{python}[noeval]
fig = px.SOORT(tabel, x="kolom", y="kolom", color="kolom",
               title="...", labels={"kolom": "Leesbare naam"})
\end{python}

\begin{itemize}
  \item de \textbf{tabel} komt eerst; daarna verwijs je naar kolommen met
        hun naam, als tekst;
  \item \code{title} is de titel boven de grafiek;
  \item \code{labels} is een dictionary die kolomnamen vertaalt naar
        leesbare aslabels, mét eenheid: ``Sprinttijd (s)'' in plaats van
        \code{sprint\_30m\_s};
  \item het resultaat is een \textbf{figuur}-object. Zet je de variabele
        als laatste regel in een cel, dan toont marimo de grafiek.
\end{itemize}

\subsection{Histogram: hoe is één variabele verdeeld?}

Een histogram verdeelt de waarden in vakjes en toont hoeveel waarden in
elk vakje vallen. \code{nbins} kiest het (ongeveer) aantal vakjes.

\begin{python}[fig=histogram]
fig = px.histogram(
    athletes, x="height_cm", nbins=15,
    title="Verdeling van de lengte van de atleten",
    labels={"height_cm": "Lengte (cm)"},
)
\end{python}

Een histogram heeft maar één kolom nodig (\code{x}); de y-as is het
aantal atleten per vakje, dat telt plotly zelf. Speel met \code{nbins}:
te weinig vakjes verbergen details, te veel maken de grafiek rommelig.

\subsection{Boxplot: verdelingen vergelijken}

Een boxplot vat een verdeling samen in vijf getallen. Zo lees je hem:

\begin{itemize}
  \item de \textbf{lijn in de doos} is de mediaan: de helft van de
        waarden ligt eronder, de helft erboven;
  \item de \textbf{doos} loopt van het eerste tot het derde kwartiel: de
        middelste helft van de waarden;
  \item de \textbf{snorharen} (\emph{whiskers}) lopen tot de laagste en
        hoogste waarde die niet uitzonderlijk ver van de doos liggen;
  \item \textbf{losse punten} zijn uitschieters.
\end{itemize}

\begin{python}[fig=boxplot]
fig = px.box(
    combined, x="position", y="sprint_30m_s",
    color="position", color_discrete_map=POSITION_COLORS,
    title="Sprinttijd over 30 m per positie",
    labels={"sprint_30m_s": "Sprinttijd (s)", "position": "Positie"},
    hover_data=["name"],
)
\end{python}

\code{hover\_data=["name"]} voegt de naam toe aan het informatievenster
dat verschijnt als je met de muis over een punt gaat. Zo zie je in het
notebook meteen wie die uitschieter is.

\subsection{Spreidingsdiagram: is er een verband?}

In een spreidingsdiagram is elke rij een punt, met de ene variabele op de
x-as en de andere op de y-as.

\begin{python}[fig=scatter]
fig = px.scatter(
    combined, x="sprint_30m_s", y="cmj_height_cm",
    color="position", color_discrete_map=POSITION_COLORS,
    title="Sprinttijd tegenover spronghoogte",
    labels={"sprint_30m_s": "Sprinttijd 30 m (s)", "cmj_height_cm": "Spronghoogte (cm)"},
    hover_data=["name", "test_moment"],
)
\end{python}

De punten vormen een brede wolk zonder duidelijke richting. Als er een
sterk verband was, zouden ze rond een schuine lijn liggen: dalend als
snellere sprinters (een lagere tijd) ook hoger springen. Hier is dat niet
zo: in deze data zegt de sprinttijd weinig over de spronghoogte. Ook
``geen duidelijk verband'' is een antwoord op de vraag van de coach. En
zie je wel een verband, onthoud dan dat dat nog geen bewijs is dat het
ene het andere veroorzaakt.

\subsection{Lijngrafiek: hoe verandert iets?}

Een lijngrafiek verbindt opeenvolgende punten. Daarvoor moet de data
\textbf{gesorteerd} zijn op de x-as, anders springt de lijn heen en weer.
Met \code{color="athlete\_id"} krijgt elke atleet een eigen lijn:

\begin{python}[fig=lijn]
fig = px.line(
    combined.sort_values("test_moment"), x="test_moment", y="sprint_30m_s",
    color="athlete_id",
    title="Sprinttijd over 30 m doorheen het seizoen, per atleet",
    labels={"sprint_30m_s": "Sprinttijd (s)", "test_moment": "Testmoment"},
)
fig.update_layout(showlegend=False)
\end{python}

Met \code{update\_layout(showlegend=False)} verberg je de legende: met 42
atleten zou die meer plaats innemen dan de grafiek zelf. Dit werkt zonder
\code{pivot\_table}, omdat \code{combined} al lang is.

% ---------------------------------------------------------------------------
\section{Meer dimensies: kleur, facetten en grootte}

Een spreidingsdiagram toont twee variabelen. Met drie extra argumenten
kun je er meer tonen:

\begin{itemize}
  \item \code{color="kolom"}: een categorie als kleur;
  \item \code{facet\_col="kolom"}: één kleine grafiek per waarde, naast
        elkaar, met dezelfde assen;
  \item \code{size="kolom"}: een getal als grootte van het punt.
\end{itemize}

\begin{python}[fig=facet]
fig = px.scatter(
    combined, x="sprint_30m_s", y="cmj_height_cm",
    color="position", color_discrete_map=POSITION_COLORS,
    facet_col="gender", size="weight_kg",
    title="Sprint tegenover sprong, per positie en geslacht",
    labels={"sprint_30m_s": "Sprinttijd 30 m (s)", "cmj_height_cm": "Spronghoogte (cm)"},
)
\end{python}

Vijf variabelen in één grafiek: sprinttijd, spronghoogte, positie,
geslacht en gewicht. Merk het derde, lege paneel \code{gender=nan} op:
bij een paar rijen ontbreekt het geslacht, en plotly maakt daar een eigen
paneel voor. Wil je dat niet, laat die rijen dan eerst weg met
\code{combined.dropna(subset=["gender"])}, en vertaal de paneeltitels
met \code{labels=\{"gender": "Geslacht"\}}. Facetten zijn vaak duidelijker dan alles in één
grafiek te proppen, omdat je de panelen naast elkaar kunt vergelijken.
Maar meer is niet altijd beter: voeg een dimensie alleen toe als ze iets
bijdraagt aan de vraag.

% ---------------------------------------------------------------------------
\section{De hartslag als tijdreeks, met zones}

Een lijngrafiek van de hartslag wordt veel leesbaarder als je de
hartslagzones als gekleurde banden op de achtergrond toont. Daarvoor
dient \code{add\_hrect}: een horizontale rechthoek tussen twee waarden op
de y-as.

\begin{python}[fig=hartslag]
ZONE_COLORS = {"Zone 1": "#cde2fb", "Zone 2": "#86b6ef", "Zone 3": "#3987e5",
               "Zone 4": "#1c5cab", "Zone 5": "#0d366b"}
one_athlete_hr = hr_zoned.loc[hr_zoned["athlete_id"] == "A001"].sort_values("timestamp")
hr_max = one_athlete_hr["hr_max"].iloc[0]

fig = px.line(
    one_athlete_hr, x="timestamp", y="heart_rate_bpm",
    title=f"Hartslag tijdens de training - {one_athlete_hr['name'].iloc[0]}",
    labels={"heart_rate_bpm": "Hartslag (bpm)", "timestamp": "Tijd"},
)
zone_edges_pct = [0, 60, 70, 80, 90, 100]
zone_names = ["Zone 1", "Zone 2", "Zone 3", "Zone 4", "Zone 5"]
for lo, hi, zone in zip(zone_edges_pct[:-1], zone_edges_pct[1:], zone_names):
    fig.add_hrect(y0=lo / 100 * hr_max, y1=hi / 100 * hr_max,
                  fillcolor=ZONE_COLORS[zone], opacity=0.35, line_width=0,
                  annotation_text=zone, annotation_position="right")
fig
\end{python}

De lus verdient wat uitleg:

\begin{itemize}
  \item \code{zone\_edges\_pct[:-1]} is de lijst zonder het laatste
        element (0, 60, 70, 80, 90): de ondergrenzen.
        \code{zone\_edges\_pct[1:]} is de lijst zonder het eerste element
        (60, 70, 80, 90, 100): de bovengrenzen.
  \item \code{zip(...)} loopt door drie lijsten tegelijk en geeft telkens
        één element van elke lijst: eerst \code{(0, 60, "Zone 1")}, dan
        \code{(60, 70, "Zone 2")}, enzovoort.
  \item \code{lo / 100 * hr\_max} zet een percentage om naar een hartslag
        in slagen per minuut.
\end{itemize}

In het notebook gebruik je \code{\_lo} in plaats van \code{lo}: in marimo
blijft een naam die met \code{\_} begint binnen de cel, zodat hij niet
botst met variabelen in andere cellen.

% ---------------------------------------------------------------------------
\section{Interactiviteit in marimo}

\subsection{UI-elementen en \code{.value}}

marimo heeft \textbf{UI-elementen}: keuzemenu's, schuifbalken,
tekstvakken. Een UI-element is een gewone Python-variabele. Het
belangrijkste kenmerk is \code{.value}: de waarde die de gebruiker op
dit moment gekozen heeft.

\begin{python}[noeval]
athlete_picker = mo.ui.dropdown(
    options=sorted(hr_zoned["name"].unique()),
    value=sorted(hr_zoned["name"].unique())[0],
    label="Atleet",
)
time_window = mo.ui.range_slider(
    start=0, stop=2400, step=60, value=(0, 2400),
    label="Tijdvenster (s)", show_value=True,
)
mo.hstack([athlete_picker, time_window])
\end{python}

\begin{itemize}
  \item \code{options} zijn de keuzes; \code{value} is de beginkeuze.
        \code{.unique()} geeft elke naam één keer, \code{sorted} zet ze
        alfabetisch.
  \item een \code{range\_slider} heeft twee handvatten; zijn
        \code{.value} is een tuple, bv.\ \code{(0, 2400)}.
  \item \code{mo.hstack([...])} zet elementen naast elkaar,
        \code{mo.vstack([...])} onder elkaar.
\end{itemize}

In een \emph{andere} cel gebruik je de gekozen waarden om de data te
filteren en de grafiek te maken:

\begin{python}[noeval]
selected_hr = hr_zoned.loc[
    (hr_zoned["name"] == athlete_picker.value)
    & hr_zoned["time_s"].between(*time_window.value)
].sort_values("timestamp")

reactive_hr_chart = px.line(selected_hr, x="timestamp", y="heart_rate_bpm",
                            title=f"Hartslag - {athlete_picker.value}")
reactive_hr_chart
\end{python}

\code{.between(a, b)} geeft een masker dat \code{True} is voor waarden
van \code{a} tot en met \code{b}. Het sterretje in
\code{between(*time\_window.value)} \emph{pakt de tuple uit}:
\code{between(*(0, 2400))} is hetzelfde als \code{between(0, 2400)}.

\begin{kernidee}[Waarom twee cellen?]
marimo voert een cel opnieuw uit als een variabele waarvan ze afhangt
verandert. Als je het keuzemenu verzet, verandert
\code{athlete\_picker.value}, en voert marimo elke cel die
\code{athlete\_picker} gebruikt opnieuw uit. De cel die het element
\emph{maakt} mag de \code{.value} niet zelf lezen: dan zou het element
bij elke keuze opnieuw gemaakt worden. Daarom: \textbf{maak een
UI-element in de ene cel, lees zijn waarde in een andere}.
\end{kernidee}

\subsection{Selecteren in de grafiek (optioneel)}

\code{mo.ui.plotly(figuur)} gaat nog een stap verder: als je in de
grafiek punten selecteert met de box- of lassotool, zitten die punten in
\code{.value}. Met \code{pd.DataFrame(selectable\_scatter.value)} krijg
je ze als tabel.

% ---------------------------------------------------------------------------
\section{Het minidashboard}

De opdracht van vandaag: een klein dashboard voor de coach, met
\textbf{drie grafieken} en \textbf{één keuzemenu} waarop minstens één
grafiek reageert. Het keuzemenu (\code{dashboard\_selector}, een keuze
uit de posities) staat al klaar in het notebook. Je maakt een lijst
\code{dashboard\_charts} met drie figuren.

Een paar tips:

\begin{itemize}
  \item \textbf{Hergebruik.} Het doel is combineren, niet nieuwe
        grafiektypes leren. Alles uit dit labo mag terugkomen.
  \item \textbf{Eén vraag per grafiek.} Schrijf eerst voor elke grafiek op
        welke vraag van de coach ze beantwoordt, en kies dan het type
        (tabel in sectie 2).
  \item \textbf{Filter met de selectie.} Filter de data met
        \code{dashboard\_selector.value} voor je de grafiek maakt, zoals
        bij de hartslaggrafiek hierboven.
  \item \textbf{Afwerking.} Elke grafiek heeft een titel en leesbare
        aslabels met eenheden, en gebruikt \code{POSITION\_COLORS} waar
        posities voorkomen.
  \item \textbf{Schikken.} Met \code{mo.vstack} en \code{mo.hstack} zet je
        het keuzemenu en de grafieken netjes onder en naast elkaar.
\end{itemize}

% ---------------------------------------------------------------------------
\section{Samenvatting}

\begin{center}
\rowcolors{2}{tint}{white}
\begin{tabularx}{\textwidth}{@{}L{7.6cm}Y@{}}
\toprule
\textbf{Code} & \textbf{Wat het doet} \\
\midrule
\code{px.histogram(df, x=, nbins=)} & verdeling van één variabele \\
\code{px.box(df, x=, y=)} & verdelingen vergelijken per groep \\
\code{px.scatter(df, x=, y=)} & verband tussen twee variabelen \\
\code{px.line(df.sort\_values(x), x=, y=)} & evolutie; eerst sorteren \\
\code{px.bar(df, x=, y=)} & waarden vergelijken; as vanaf 0 \\
\code{title=}, \code{labels=\{...\}}, \code{hover\_data=[...]} & titel, leesbare aslabels, extra info bij hover \\
\code{color=}, \code{color\_discrete\_map=} & categorie als kleur, met vaste kleuren \\
\code{facet\_col=}, \code{size=} & panelen per waarde, getal als puntgrootte \\
\code{fig.update\_yaxes(rangemode="tozero")} & y-as laten beginnen bij 0 \\
\code{fig.add\_hrect(y0=, y1=, ...)} & horizontale band, bv.\ een hartslagzone \\
\code{mo.ui.dropdown}, \code{mo.ui.range\_slider} & keuzemenu, schuifbalk met twee handvatten \\
\code{element.value} & de huidige keuze van de gebruiker \\
\code{mo.hstack}, \code{mo.vstack} & elementen naast of onder elkaar zetten \\
\bottomrule
\end{tabularx}
\end{center}

\subsection*{Veelgemaakte fouten}

\begin{itemize}
  \item Een staafdiagram met een y-as die niet bij 0 begint.
  \item Geen titel of geen eenheden op de assen.
  \item Een lijngrafiek op data die niet gesorteerd is op de x-as.
  \item Een UI-element maken en zijn \code{.value} in dezelfde cel lezen.
  \item Elke grafiek eigen kleuren laten kiezen, zodat ``Forward'' in de
        ene grafiek blauw is en in de andere oranje.
\end{itemize}

\subsection*{Begrippen}

\begin{description}[leftmargin=3.2cm, labelwidth=3.0cm, font=\normalfont\bfseries]
  \item[figuur] het object dat een \code{px}-functie teruggeeft; marimo toont het als grafiek.
  \item[mediaan] middelste waarde: de helft ligt eronder, de helft erboven.
  \item[kwartiel] grens van de laagste of hoogste 25\,\% van de waarden.
  \item[uitschieter] waarde die ver van de andere ligt.
  \item[facet] klein paneel met dezelfde assen, één per waarde van een kolom.
  \item[UI-element] keuzemenu, schuifbalk of ander element waarmee de gebruiker iets kiest.
\end{description}

\end{document}
